\section{Why Action Enters the Update in Closed-Loop Tracking}\label{sec:necessity}

The next step is to justify the minimal closed-loop architecture used
throughout the paper. The finite-memory law studied later concerns tracking
error once a representation is already in place, but in action-effective
systems there is a prior question: what must an update rule remember in order
to remain coherent under feedback? In open-loop settings, combining past state
and new observations is enough. In closed-loop settings, however, the system's
own actions alter future observations. The update therefore needs to keep track
of what the system did, not just of what it saw.

\subsection{External Characterization: Drift-Coherence}

\begin{definition}[Drift-Coherence]\label{def:driftcoherent}
An information system $\IS$ is \emph{drift-coherent} if there exist
$\beta \in \mathcal{KL}$ and $\gamma \in \mathcal{K}$ such that for all
$t \geq 0$ and all environmental inputs with
$\mathrm{ess\,sup}_{s \leq t}\|u(s)\| \leq M < \infty$:
\[
  W_2\!\left(\Pot(S,t),\,\Pot^*(S,t)\right)
  \;\leq\;
  \beta\!\left(W_2\!\left(\Pot(S,0),\,\Pot^*(S,0)\right),\,t\right)
  \;+\;
  \gamma(M).
\]
The system maintains bounded tracking error under bounded drift.
\end{definition}

\begin{definition}[Action-Effectiveness]\label{def:actioneffective}
The actuation operator $\Act_t$ is \emph{action-effective} in environment
$\Env$ if there exist distinct outputs $o, o' \in \mathcal{O}$ and a state
$s \in S$ such that
\[
  D_{\mathrm{TV}}\!\left(P(\cdot \mid s_t = s,\, o)
  \;\Big\|\;
  P(\cdot \mid s_t = s,\, o')\right) \;>\; 0.
\]
If this never occurs, the system is \emph{passive}: its actions do not affect
the environmental transition law.
\end{definition}

\subsection{Belief Update With Feedback}

\begin{lemma}[Belief Update Lemma]\label{lem:beliefupdate}
Let $\IS$ be a drift-coherent, action-effective information system. Any update
map $U$ whose repeated application preserves drift-coherence in this class must
depend on three arguments:
\begin{enumerate}[label=\textnormal{(\roman*)}]
  \item the current internal state $p = \Pot(S,t)$,
  \item the most recent action $o = \Act_t(\Pot)$, and
  \item the incoming environmental signal $s_{t+1}$.
\end{enumerate}
Each argument cannot be dropped without leaving the class of drift-coherent
updates considered here.
\end{lemma}

\begin{proof}
If $U$ omits $p$, the update becomes memoryless and discards the sufficient
summary of the past; under temporally structured drift this prevents bounded
tracking error. If $U$ omits $o$, then in an action-effective environment the
update cannot distinguish exogenous drift from changes induced by the system's
own actuation. This is precisely the credit-assignment problem that appears in
POMDP filtering and belief update semantics
\citep{kaelbling1998,KatsunaMendelzon1992,Bonanno2024}. If $U$ omits
$s_{t+1}$, no new environmental evidence enters the state estimate and the gap
between $\Pot$ and $\Pot^*$ grows with drift. In each case, the
drift-coherence bound fails.
\end{proof}

\begin{remark}[Passive observer systems]\label{rem:passive}
When actions do not influence transitions, the update can drop the actuation
argument and the system is genuinely dyadic. The three-role decomposition is
therefore motivated only for action-effective settings.
\end{remark}

\begin{corollary}[Ternary Update Structure]\label{cor:ternary}
For action-effective tracking systems, any drift-coherent update map has the
form
\[
  U\colon \Info \times \mathcal{O} \times S \to \Info.
\]
Any ostensibly smaller description either hides one of these arguments inside a
state variable or fails drift-coherence.
\end{corollary}

\begin{proof}
This is immediate from \cref{lem:beliefupdate}.
\end{proof}

\begin{theorem}[Minimal Functional Roles]\label{thm:irreducibility}
Let $\IS$ be a drift-coherent, action-effective information system. Any
functional decomposition that preserves \cref{cor:ternary} in this class must
realize three roles: a stored state, an action channel, and an adaptation rule
that maps the state, action, and new observation to an updated state. The tuple
$(\Pot,\Act,\Conv)$ is the explicit decomposition used here.
\end{theorem}

\begin{proof}
By \cref{cor:ternary}, the update requires all three arguments
$(p,o,s_{t+1})$. A decomposition that omits any one of the corresponding roles
must either re-encode it implicitly inside another component or lose the update
structure required for drift-coherence. The explicit realization used here is
$\Pot$ for stored state, $\Act$ for the action channel, and $\Conv$ for the
update map.
\end{proof}

This necessity result does not yet tell us when such a three-role system is
stable or diagnostically meaningful. It only shows that once actions affect
future observations, a dyadic state-observation update is too small.
The next section supplies the corresponding stability template under which the
three roles can be scored.
